\subsection{Градиент и гессиан логистической регрессии}

Для логистической регрессии с $b_i\in\{-1,1\}$ и признаками $a_i\in\R^d$
\[
    f(w)=
    \frac{1}{m}\sum_{i=1}^{m}
    \log(1+\exp(-b_i\inner{a_i}{w}))
    +
    \frac{\lambda}{2}\norm{w}_2^2.
\]
Обозначим
\[
    z_i=b_i\inner{a_i}{w},
    \qquad
    \sigma_i=\frac{1}{1+\exp(z_i)}.
\]
Тогда градиент равен
\[
    \nabla f(w)
    =
    -\frac{1}{m}\sum_{i=1}^{m}b_i\sigma_i a_i
    +
    \lambda w.
\]
Гессиан имеет вид
\[
    \nabla^2 f(w)
    =
    \frac{1}{m}
    \sum_{i=1}^{m}
    \sigma_i(1-\sigma_i)a_i a_i^T
    +
    \lambda I.
\]
Даже если матрица данных разреженная, явное построение гессиана требует памяти порядка
$d^2$.
Поэтому используются разности градиентов как практический способ получить информацию
о локальной кривизне без хранения матрицы.

\subsection{Линейный оракул на $l_1$-шаре}

Для ограничения
\[
    \norm{w}_1\leq R
\]
линейная подзадача FW имеет вид
\[
    \min_{\norm{s}_1\leq R}\inner{g}{s}.
\]
Оптимум достигается в вершине $l_1$-шара.
Пусть
\[
    j\in\argmax_{\ell}|g_\ell|.
\]
Тогда
\[
    s=-R\,\operatorname{sign}(g_j)e_j.
\]
Это очень дешевая операция: нужно найти координату максимального по модулю градиента.
В этом смысле FW хорошо согласуется с разреженными моделями.
Если стартовать из нуля или из вершины $l_1$-шара, итераты часто остаются разреженными
в виде выпуклых комбинаций небольшого числа вершин.

\subsection{Роль радиуса $R$}

Радиус $R$ задает силу ограничения.
Если решение без ограничений имеет норму намного больше $R$, то оптимум ограниченной задачи
лежит на границе $l_1$-шара.
В этом режиме FW часто страдает от зигзагообразного поведения: линейный оракул выбирает
то одну, то другую вершину, а линейный поиск делает небольшие шаги вдоль ребер допустимого
множества.
Именно здесь CFW и NFW наиболее полезны.

Если решение без ограничений лежит внутри шара, то ограничение неактивно.
Тогда FW фактически пытается решить гладкую задачу с искусственным компактным множеством.
Линейный оракул по вершинам шара может быть менее согласован с естественной геометрией
минимума, и различия между FW с линейным поиском и CFW/NFW становятся меньше.

Это проверяется через четыре набора радиусов:
\[
\begin{array}{c|c}
\text{датасет} & R \\
\hline
\text{Mushrooms} & 5,\;30,\;60,\;150\\
\text{MNIST} & 0.10,\;0.80,\;1.50,\;2.50\\
\text{MNIST с }l_2 & 0.01,\;0.10,\;0.20,\;0.50
\end{array}
\]
Эти значения специально выбираются так, чтобы покрыть разные положения оптимума относительно
границы допустимого множества.

\subsection{Линейный поиск в ML-экспериментах}

Для всех модификаций, кроме FW Vanilla, используется линейный поиск:
\[
    \gamma_k=\argmin_{\gamma\in[0,1]}f(w_k+\gamma d_k).
\]
Для логистической регрессии одномерная функция гладкая и выпуклая.
Ее производная равна
\[
    \varphi'(\gamma)
    =
    \inner{\nabla f(w_k+\gamma d_k)}{d_k}.
\]
Поэтому можно применять стандартный одномерный поиск.
FW Vanilla использует теоретический шаг
\[
    \gamma_k=\frac{2}{k+2}.
\]
В таблицах видно, что FW Vanilla часто проигрывает FW с линейным поиском.
Однако основной вопрос состоит не в выборе шага, а в выборе направления: CFW и NFW сравниваются
с сильной базовой версией FW ls.

\subsection{Особенности результатов на MNIST}

MNIST имеет более сложную и более высокоразмерную структуру признаков.
Линейный оракул на $l_1$-шаре выбирает одну координату, но полезная информация может быть
распределена по многим пикселям.
В такой ситуации вершины допустимого множества становятся плотнее с точки зрения возможных
направлений, и разница между обычным FW-направлением и сопряженной комбинацией уменьшается.
Поэтому на MNIST CFW и NFW часто выигрывают у FW по накопленному зазору, но не дают такого
резкого отрыва, как на Mushrooms.

\subsection{Итоги раздела}

Рассмотренный подход отвечает на вопрос, можно ли перенести транспортные сопряженные FW-методы
в машинное обучение.
Прямое препятствие состоит в стоимости второго порядка: явное построение гессиана и точные
произведения гессиана на направления слишком дороги для типичных ML-задач.
Главный технический шаг состоит в аппроксимации этих произведений конечными разностями
градиентов.
После этого CFW и NFW сохраняют почти ту же итерационную структуру, что и FW.

Теоретический результат показывает, что CFW не теряет базовый порядок сходимости FW.
Экспериментальный результат показывает, что в задачах, где ограничение активно и вызывает
зигзагообразное движение, сопряженные направления дают существенное практическое ускорение.
